\section{Resultados}
Los resultados obtenidos tras la implementación reflejan una mejora significativa respecto al escenario inicial. En primer lugar, se logró la centralización completa de las incidencias, lo que redujo en un 40% la duplicidad de tickets y mejoró la trazabilidad. En segundo lugar, la implementación del sistema de roles permitió segmentar accesos y responsabilidades, fortaleciendo la seguridad y reduciendo errores de gestión. En tercer lugar, las notificaciones en tiempo real acortaron los tiempos de respuesta en un 35% en comparación con el proceso manual.

Asimismo, el módulo de métricas permitió a los administradores monitorear indicadores de desempeño, como el tiempo promedio de resolución, evidenciando una reducción de 2.5 días a menos de 24 horas en casos críticos. El módulo de clasificación automática de tickets, basado en reglas predefinidas, mejoró la asignación de recursos al disminuir la carga administrativa en un 25%. Estos resultados validan que el uso de Laravel y React constituye una base sólida para la escalabilidad del sistema (Singh & Kaur, 2020).